\section{Statement of the Problem}
Hospitals must provide essential medical services during and after earthquakes. Structural vibration and foundation flexibility influence seismic displacement and drift, yet a fixed-base model does not represent the deformability of all supporting soils \cite{guragain2004,subedi2022}. This motivates an assessment of how foundation-support assumptions alter the seismic response quantities reported for a real hospital in Pokhara.

The case-study architectural documentation identifies six structural blocks within the G+4 hospital. The central question is how each block's reported modal response changes between fixed-base, hard-, medium-, and soft-soil cases, and how the associated displacement, drift and fragility charts describe differences among those support conditions. A coherent comparison requires a consistent block identifier, structural response quantity, support case, and performance-state definition.
